\clearpage
\fancyhead[L]{\small KONSTANTEN DES BIFILARMAGNETOMETERS.}
\SourceLeaf{410}\pdfbookmark[0]{Zur Bestimmung der Konstanten des Bifilarmagnetometers.}{bifilarconstanten-title}
\begin{center}
{\Large\bfseries ZUR BESTIMMUNG\par\medskip DER KONSTANTEN DES BIFILARMAGNETOMETERS.}
\par\bigskip
Resultate aus den Beobachtungen des magnetischen Vereins. 1840. I.
\end{center}
\renewcommand{\Article}[1]{\par\medskip\Needspace{4\baselineskip}\pdfbookmark[1]{#1.}{bifilarconstanten-art-#1}\centerline{\textbf{#1.}}\nobreak\smallskip}%
\Article{1}
Zum richtigen und sicheren Gebrauch des Bifilarmagnetometers muss man die Zahlenwerte gewisser Größen kennen. Sie beziehen sich auf Verhältnisse der Teile des Apparats, die unter bestimmten Bedingungen als konstant betrachtet werden können. Von diesen wesentlichen Größen hängen die Gleichgewichtslagen und Schwingungszeiten ab, die bei verschiedenen Stellungen der beweglichen Teile zu beobachten sind. Es handelt sich um vier Größen.

1) Die Stellung, die der Index der Spiegelalhidade haben muss, damit die Normale zum Spiegel und die optische Achse des Beobachtungsfernrohrs in derselben Vertikalebene liegen, wenn die beiden Aufhängungsdrähte in einer Vertikalebene liegen. Diese Stellung soll mit $P$ bezeichnet werden, wobei die reflektierende Fläche des Spiegels dem Fernrohr zugekehrt ist.

2) Die Stellung, die man bei eben dieser Lage der Aufhängungsdrähte dem Index des Schiffchens geben muss, damit die magnetische Achse des Magnetstabs in ihrer natürlichen Lage im magnetischen Meridian liegt. Ich bezeichne diese Stellung mit $Q$.

Es versteht sich, dass man, wenn jede der beiden Alhidaden mehr als einen Index hat, immer einen davon nach Belieben als Hauptindex wählen muss.

3) Das Verhältnis der magnetischen Richtkraft zu der aus der Art der Aufhängung entstehenden Richtkraft. Diese letztere kann statische Richtkraft heißen. Das Verhältnis soll durch $R:1$ ausgedrückt werden.

\SourceLeaf{411}4) Die statische Schwingungsdauer des Apparats, also diejenige, die sich allein aus der Art der Aufhängung ergeben würde, ohne Einwirkung des Erdmagnetismus auf den Magnetstab. Ich bezeichne das Quadrat dieser Schwingungsdauer mit $S$.

Daraus geht hervor, dass $\frac{S}{R}$ das Quadrat der rein magnetischen Schwingungsdauer ausdrückt, also derjenigen, die sich bei der Aufhängung des Apparats an einem einzigen Faden ohne Torsion ergeben würde.

\Article{2}
Zunächst ist herzuleiten, wie die unmittelbaren Ablesungen am Bifilarmagnetometer mit der Stellung der beiden Alhidaden und diesen vier Größen zusammenhängen.

Steht die Spiegelalhidade auf $A$ und die Alhidade des Schiffchens auf $B$, so bezeichne $t$ die Schwingungsdauer und $p$ den in Bogenteile umgerechneten Abstand des zur Gleichgewichtslage gehörenden Skalenteils von dem Punkt der Skala, der mit der optischen Achse des Beobachtungsfernrohrs in derselben Vertikalebene liegt. Dieser Punkt ist durch den von der Mitte des Objektivs herabhängenden Lotfaden gekennzeichnet. Um die Richtungen eindeutig festzulegen, nehme ich an, dass die Teilungen am Kreis und an der Skala von links nach rechts verlaufen, und rechne $p$ positiv, wenn die am Fadenkreuz des Fernrohrs beobachtete Zahl größer ist als die Zahl am Lotfaden. In dieser Gleichgewichtslage ist das Bifilarmagnetometer also um $A-P-p$ rückwärts, das heißt von rechts nach links, gegenüber der Lage gedreht, in der die Aufhängungsdrähte parallel waren. Anders gesagt ist $A-P-p$ der Winkel zwischen der Geraden durch die beiden unteren Enden der Aufhängungsdrähte und einer Parallelen zu der Geraden, die ihre beiden oberen Enden verbindet. Das durch die Aufhängung hervorgebrachte Drehmoment ist dem Sinus dieses Winkels zwar nicht vollkommen, aber für die praktische Anwendung hinreichend genau proportional. Wir setzen es $=D\sin(A-P-p)$, wobei $D$ die statische Richtkraft bezeichnet. Positive Werte des Drehmoments beziehen sich auf eine Drehung von links nach rechts.

In derselben Lage des Apparats bildet die magnetische Achse des Magnetstabs mit dem magnetischen Meridian den von rechts nach links ge%
\expandafter\def\csname HistoricalNoteText416\endcsname{An der Nadel eines Kompasses auf einer 1709 angefertigten Sonnenuhr der hiesigen Sternwarte ließ sich der Magnetismus 1841 durch erneutes Bestreichen bis zur Sättigung nur auf das Dreifache steigern, bei einer anderen von 1603 nur auf das Fünffache. Unter der sehr wahrscheinlichen Voraussetzung, dass beide seit ihrer Herstellung nie neu bestrichen worden waren, und wenn man zugleich annimmt, dass sie ursprünglich bis zur Sättigung magnetisiert waren und die Kraft allmählich in geometrischer Folge abgenommen hat, beträgt der jährliche Verlust bei der ersten $\frac{1}{120}$, bei der zweiten $\frac{1}{150}$. Er ist noch kleiner, falls die ursprüngliche Magnetisierung die Sättigung nicht erreicht hatte.}%
\SourceLeaf{412}zählten Winkel $A-P-p-B+Q$. Das Drehmoment, das aus der Einwirkung des Erdmagnetismus auf den Magnetstab entsteht und von links nach rechts positiv gerechnet wird, ist $=RD\sin(A-P-p-B+Q)$. Damit haben wir die Gleichung (1)
\[
0=\sin(A-P-p)+R\sin(A-P-p-B+Q)
\]
Dreht man den ganzen Apparat aus der Gleichgewichtslage um den Winkel $z$ von rechts nach links, so wirkt im entgegengesetzten Sinn das Drehmoment
\[
D\sin(z+A-P-p)+DR\sin(z+A-P-p-B+Q)
\]
Entwickelt man die beiden Sinusausdrücke und berücksichtigt Gleichung (1), so ergibt sich
\[
D\sin z(\cos(A-P-p)+R\cos(A-P-p-B+Q))
\]
Dieser Ausdruck ist also dem Sinus von $z$ proportional. Man muss daher
\[
D(\cos(A-P-p)+R\cos(A-P-p-B+Q))
\]
als die Richtkraft ansehen, die sich aus der Verbindung der statischen und der magnetischen Richtkraft ergibt, und wir haben daher (2)
\[
\frac{S}{tt}=\cos(A-P-p)+R\cos(A-P-p-B+Q)
\]
Quadriert man in den beiden Gleichungen (1), (2) jeweils beide Seiten und addiert sie, so findet man (3)
\[
\frac{SS}{t^4}=1+2R\cos(Q-B)+RR
\]
Bezeichnet man mit $e$ die Basis der hyperbolischen Logarithmen und mit $i$ die imaginäre Einheit $\BareRoot{-1}$, so lassen sich die beiden Gleichungen (1), (2) bequem zu einer Gleichung zusammenfassen
\[
\frac{S}{tt}=e^{i(A-P-p)}+Re^{i(A-P-p-B+Q)}
\]
oder noch einfacher zur folgenden Gleichung (4)
\[
1=\frac{S}{tt}e^{i(P+p-A)}-Re^{i(Q-B)}
\]
\SourceLeaf{413}die die beiden Gleichungen
\[
\begin{gathered}
1=\frac{S}{tt}\cos(P+p-A)-R\cos(Q-B)\\
\frac{S}{tt}\sin(P+p-A)=R\sin(Q-B)
\end{gathered}
\]
umfasst.

Für die natürliche Lage, in der $Q=B$ gilt, ergibt sich
\[
tt=\frac{S}{1+R}
\]
für die umgekehrte dagegen, in der $Q=B+180^0$ gilt,
\[
tt=\frac{S}{1-R}
\]

\Article{3}
Die transversale Stellung im engeren Sinn erfordert, dass
\[
A-P-p-B+Q=\pm90^0
\]
gilt, wobei das obere Zeichen für den Fall steht, dass der Nordpol des Magnetstabs westlich des magnetischen Meridians liegen soll, das untere für die östliche Lage. Nach (1) ergibt sich also
\[
\sin(A-P-p)=\mp R
\]
Bezeichnet man demnach mit $\varphi$ den spitzen Winkel, dessen Sinus $=R$ ist, so gilt für die westliche Stellung des Nordpols
\[
A=P+p-\varphi,\quad B=Q-\varphi-90^0
\]
für die östliche dagegen
\[
A=P+p+\varphi,\quad B=Q+\varphi+90^0
\]
Damit die Gleichgewichtslage dem durch den Lotfaden bezeichneten Skalenpunkt selbst entspricht, muss man die Spiegelalhidade im ersten Fall auf $P-\varphi$ und im zweiten auf $P+\varphi$ einstellen.

Für die zur Transversalstellung gehörende Schwingungsdauer ergibt Formel (2) sofort
\SourceLeaf{414}\[
\frac{S}{tt}=\cos\varphi
\]
oder
\[
tt=\frac{S}{\BareRoot{(1-RR)}}
\]
Die Schwingungsdauer in der Transversalstellung ist demnach das geometrische Mittel der Schwingungszeiten in der natürlichen und der umgekehrten Stellung.

\Article{4}
Um klar zu erkennen, inwieweit die Größen veränderlich sind, müssen wir sie auf ihren Ursprung zurückführen.

Die Winkel $P$ und $Q$ setzen sich jeweils aus drei Teilen zusammen. Der erste Teil von $P$ ist der Winkel zwischen dem Radius zum Nullpunkt des Kreises und der Geraden, die die beiden unteren Enden der Aufhängungsdrähte verbindet. Der zweite ist der Winkel zwischen der Geraden, die die beiden oberen Enden der Aufhängungsdrähte verbindet, und der optischen Achse des Fernrohrs (oder vielmehr zwischen den Projektionen dieser Geraden auf eine Horizontalebene, was auch bei allen anderen Winkelschenkeln stillschweigend vorausgesetzt wird, die nicht selbst horizontal sind oder sich nicht unmittelbar schneiden). Der dritte ist der Winkel zwischen der Normalen zum Spiegel und dem Radius zum Hauptindex der Spiegelalhidade.

Der erste Teil von $Q$ stimmt mit dem ersten Teil von $P$ überein. Der zweite ist der Winkel zwischen der Geraden, die die beiden oberen Enden der Aufhängungsdrähte verbindet, und dem magnetischen Meridian. Der dritte ist der Winkel zwischen der magnetischen Achse des im Schiffchen liegenden Magnetstabs und dem Radius zum Hauptindex der Alhidade des Schiffchens.

Alle diese fünf Winkel sind von links nach rechts zu zählen. Aus dieser Zerlegung geht hervor, dass $P$ ganz unveränderlich bleibt, solange die Aufhängung des Instruments, die Verbindung des Spiegels mit seiner Alhidade und die Stellung des Beobachtungsfernrohrs unverrückt bleiben. Wegen seines zweiten Bestandteils erfährt $Q$ aber genau dieselben Änderungen wie die magnetische Deklination, wobei Änderungen von links nach rechts als positiv gelten.

\SourceLeaf{415}Die statische Richtkraft wird durch die Formel
\[
D=\frac{fgG}{4h}
\]
ausgedrückt, wobei $G$ das Gewicht des Apparats (also die mit der Schwerkraft multiplizierte Masse), $f$ den Abstand der Aufhängungsdrähte an den unteren Enden, $g$ ihren Abstand an den oberen Enden und $h$ die Höhe der oberen Befestigung über der unteren bezeichnet. Das gilt zumindest, wenn man die geringe Vergrößerung außer Betracht lässt, die diese Kraft noch durch den Widerstand der einzelnen Aufhängungsdrähte gegen Torsion erfährt. Das kann man hier ohne Weiteres tun, da es zunächst nur um die Veränderlichkeit der gesamten Kraft geht. Bezeichnet man noch das Trägheitsmoment bezüglich der vertikalen Drehachse mit $K$, so ergibt sich mit $\pi$ in der üblichen Bedeutung
\[
S=\frac{\pi\pi K}{D}=\frac{4\pi\pi hK}{fgG}
\]
Es ist nun klar, dass sich die einzelnen Faktoren $f,g,h,K$ infolge von Temperaturänderungen ändern. Diese Änderungen sind teils an sich sehr gering, teils gleichen sie sich, wie weiter unten gezeigt wird, im Wert von $S$ fast vollständig aus. Der Unterschied, der sich aus ungleichem Gewichtsverlust bei ungleicher Luftdichte ergibt, kann als ganz unmerklich gelten.

Die magnetische Richtkraft ist $=TM$, wenn $T$ die Intensität der horizontalen Komponente des Erdmagnetismus und $M$ das magnetische Moment des Magnetstabs bezeichnet. Wir haben demnach
\[
R=\frac{TM}{D}=\frac{4TMh}{fgG}=\frac{STM}{\pi\pi K}
\]

Die Veränderlichkeit von $R$ hat also drei Ursachen.

Erstens die ständige Veränderlichkeit von $T$. Zweitens die Veränderlichkeit der Temperatur, die sowohl die linearen Größen $f,g,h$ als auch den Stabmagnetismus $M$ beeinflusst. Drittens die Veränderlichkeit von $M$ unabhängig von der jeweiligen Temperatur.

Über die dritte Ursache wissen wir bisher noch recht wenig. Bei den im 2. Band der \textit{Resultate} mitgeteilten Versuchen von Professor \textsc{Weber} wurde der durch künstliche Erwärmung erlittene Verlust durch die anschließende Abkühlung niemals vollständig ausgeglichen.
\SourceLeaf{416}Nach Wiederherstellung der anfänglichen Temperatur blieb ein beträchtlicher dauerhafter Verlust. Andererseits zeigt die Erfahrung, dass Magnetnadeln trotz der täglichen und jährlichen Temperaturwechsel ohne erneutes Bestreichen viele Jahre lang einen beträchtlichen Grad an Magnetismus behalten. Daraus muss man auf einen äußerst langsamen fortschreitenden Verlust schließen\HistoricalNote{416}. Es wäre sehr wichtig, die Bedingungen genau zu kennen, unter denen der Temperaturwechsel den kleinstmöglichen dauerhaften Kraftverlust bewirkt. Außer auf die Beschaffenheit und Härtung des Stahls und eine starke anfängliche Magnetisierung kommt es wahrscheinlich vor allem darauf an, dass seit dieser Magnetisierung eine gewisse Zeit vergangen ist, dass die Temperaturänderungen bestimmte Grenzen einhalten und dass sie immer nur sehr langsam und allmählich erfolgen. Unter diesen Bedingungen dürfte man den magnetischen Zustand eines Magnetstabs (wenn wir damit sein auf eine bestimmte Normaltemperatur reduziertes magnetisches Moment bezeichnen) für eine mäßige Zeit, etwa einige Tage, als konstant betrachten. Wird nach längerer Zeit eine deutliche Abnahme festgestellt, darf man für die Zwischenzeit eine stetige Abnahme in geometrischer Folge zugrunde legen. Die Ausführung des einfallsreichen Vorschlags von Professor \textsc{Weber} in dem weiter unten folgenden Aufsatz scheint besonders geeignet, Licht in diese interessante Frage zu bringen.

\Article{5}
Damit die Aufgabe, die Zahlenwerte der Größen eines Bifilarmagnetometers durch Versuche zu bestimmen, eine genaue Bedeutung erhält, verstehen wir unter den gesuchten Werten der veränderlichen Größen diejenigen, die sich auf eine bestimmte Deklination, eine bestimmte horizontale Intensität, eine bestimmte Temperatur und den magnetischen Zustand des
\SourceLeaf{417}Magnetstabs zur Zeit dieser Versuche beziehen. Änderungen, die der Stab nach längerer Zwischenzeit erfahren kann, kommen dabei also gar nicht in Betracht. Wir bezeichnen diese Normalwerte der veränderlichen Größen mit $Q^0,R^0,S^0$ (da $P$ schon für sich konstant ist) und setzen allgemein
\[
Q=Q^0+q,\quad R=rR^0,\quad S=sS^0
\]
Entsprechend mögen $f^0,g^0,h^0,K^0,T^0,M^0$ die Normalwerte der veränderlichen Größen $f,g,h,K,T,M$ bezeichnen. Damit haben wir sofort
\[
s=\frac{f^0g^0hK}{fgh^0K^0},\qquad r=\frac{f^0g^0hMT}{fgh^0M^0T^0}
\]

Um bei der Bestimmung der Größen die Änderungen der Richtung und Stärke der erdmagnetischen Kraft während der dazu nötigen Vorgänge berücksichtigen zu können, muss natürlich ein Hilfsapparat zur Verfügung stehen, am besten ein Unifilarmagnetometer, an dem man Schwingungsdauer und Stellung zugleich beobachtet. Dieses Hilfsmagnetometer dient auch dazu, die zu wählende Normaldeklination und Normalintensität nachweisbar zu machen. Dazu lässt man zunächst die erste einem bestimmten Skalenpunkt und die zweite einer bestimmten Schwingungsdauer bei der Normaltemperatur entsprechen. Damit hat man es auch in der Hand, beide Normalgrößen nach bekannten Methoden in absoluten Maßen auszudrücken. Danach ist $q$ unmittelbar die in Bogenteile umgerechnete Differenz der am Hilfsmagnetometer beobachteten Stellung gegenüber der Normalstellung. Bezeichnet man ferner die Größen, die am Bifilarmagnetometer $M,K,t$ heißen, am Hilfsmagnetometer mit $m,k,\theta$ und ihre Normalwerte mit $m^0,k^0,\theta^0$, so ergibt sich
\[
\frac{\theta\theta}{\theta^0\theta^0}=\frac{km^0T^0}{k^0mT}
\]
und folglich
\[
r=\frac{f^0g^0hkm^0M\theta^0\theta^0}{fgh^0k^0mM^0\theta\theta}=\frac{skK^0m^0M\theta^0\theta^0}{k^0KmM^0\theta\theta}
\]

Von den sieben Faktoren $\frac{f^0}{f},\frac{g^0}{g},\frac{h}{h^0},\frac{k}{k^0},\frac{K^0}{K},\frac{m^0}{m},\frac{M}{M^0}$ in den Ausdrücken für $s$ und $r$ sind die ersten fünf aus der temperaturbedingten Ausdehnung der betreffenden Metalle zu berechnen, die beiden letzten dagegen nach der besten verfügbaren Kenntnis vom Einfluss der Tem\SourceLeaf{418}peratur auf den Stabmagnetismus. Dabei wird das, was wir den magnetischen Zustand genannt haben, bei beiden Magnetstäben während der hier betrachteten Vorgänge als konstant angesehen. Zu dieser Rechnung fügen wir noch einige Ausführungen hinzu.

Wir wählen den Gefrierpunkt als Normaltemperatur und bezeichnen mit $c$ und $c'$ die Temperatur im Kasten des Bifilar- und des Hilfsmagnetometers, mit $c''$ die Temperatur an der oberen Befestigung der Aufhängungsdrähte des ersteren. Für eine Temperaturzunahme um ein Grad bezeichnen wir ferner die Ausdehnung des Stahls mit $\alpha$, die des Messings mit $\beta$ und die Abnahme des Stabmagnetismus für die Stäbe der beiden Apparate mit $\gamma$ und $\gamma'$. Da die Änderung des Trägheitsmoments der beiden Apparate zum weitaus größten Teil von der Ausdehnung der Magnetstäbe selbst herrührt, kann man ohne Bedenken
\[
\frac{K}{K^0}=(1+\alpha c)^2,\qquad \frac{k}{k^0}=(1+\alpha c')^2
\]
setzen. Für die Ausdehnung der Aufhängungsdrähte kann man, wenn sie wie beim hiesigen Apparat aus Stahl bestehen, denselben Koeffizienten $\alpha$ beibehalten und für ihre Temperatur $\tfrac{1}{2}(c+c'')$ annehmen, sodass
\[
\frac{h}{h^0}=1+\tfrac{1}{2}\alpha(c+c'')
\]
gilt. Wir haben somit (1)
\[
s=\frac{(1+\alpha c)^2(1+\tfrac{1}{2}\alpha(c+c''))}{(1+\beta c)(1+\beta c'')}
\]
wofür man auch mit hinreichender Genauigkeit
\[
s=1+(3\alpha-2\beta)c-(\beta-\tfrac{1}{2}\alpha)(c''-c)
\]
setzen kann. Da nun der Erfahrung zufolge sehr nahe $\beta=\tfrac{3}{2}\alpha$ gilt, ergibt sich in sehr guter Annäherung (2)
\[
s=1-\alpha(c''-c)
\]
d i. die Änderung der Größe $S$ hängt nur vom Unterschied zwischen der unteren und der oberen Temperatur ab, sodass $S$ in der Regel als ganz konstant angesehen werden kann.

Wir haben ferner (3)
\[
r=\frac{s\theta^0\theta^0}{\theta\theta}.\frac{1-\gamma c}{1-\gamma'c'}.\left(\frac{1+\alpha c'}{1+\alpha c}\right)^2
\]
\SourceLeaf{419}oder, wenn sich die Temperaturänderungen auf beide Stäbe gleich auswirken, das heißt, wenn $\gamma'=\gamma$ gilt, mit hinreichender Genauigkeit,
\[
r=\frac{s\theta^0\theta^0}{\theta\theta}(1+(\gamma+2\alpha)(c'-c))
\]
oder nach (2) ebenso genau (4)
\[
r=\frac{\theta^0\theta^0}{\theta\theta}(1+(\gamma+2\alpha)(c'-c)-\alpha(c''-c))
\]

Schließlich ist noch zu beachten, dass man aus den Vergleichsbeobachtungen am Unifilarmagnetometer keinen für einen bestimmten Augenblick geltenden Wert von $\frac{T}{T^0}$ ableiten kann, sondern nur den Mittelwert über die ganze Zeit, die die Schwingungsbeobachtungen umfassen. Es versteht sich also von selbst, dass sich auch alle anderen Größen, mit denen diese Schwingungsbeobachtungen als gleichzeitig unmittelbar oder mittelbar verbunden werden sollen, als Mittelwerte auf denselben Zeitraum beziehen müssen.

\Article{6}
Die einfachste Art, die vier Größen zu bestimmen, ist die folgende.

Bei beliebiger Stellung des Schiffchens legt man statt des Magnetstabs einen nichtmagnetischen Stab von ungefähr gleichem Gewicht hinein. Den Spiegel stellt man so ein, dass in der Gleichgewichtslage das Bild irgendeines Skalenpunkts auf dem Fadenkreuz des Beobachtungsfernrohrs erscheint. Nimmt man $A$ und $p$ in der obigen Bedeutung, so ergibt sich $P=A-p$. Um das Ergebnis von einer sehr genauen Kenntnis des Werts der Skalenteile oder von einer sehr genauen Umrechnung in Bogenteile unabhängiger zu machen, kann man den Vorgang, wenn $p$ beim ersten Mal noch sehr groß war, mit einer neuen, der gesuchten sehr nahe kommenden Spiegelstellung wiederholen. Am besten eignet sich dafür ein mit Blei beschwerter Holzstab. Das ungefähr gleiche Gewicht ist nötig, damit eine geringe Torsion, die die Aufhängungsdrähte für sich genommen in der Gleichgewichtslage des Ganzen möglicherweise haben, ohne Wirkung bleibt.

Ohne die Stellung des Spiegels weiter zu verändern, legt man statt der bisherigen Belastung den Magnetstab in das Schiffchen, das dann so eingestellt
\SourceLeaf{420}werden soll, dass der Ruhestellung derselbe Skalenpunkt entspricht wie zuletzt bei der nichtmagnetischen Belastung. Dazu ermittelt man durch Versuche zwei verschiedene Stellungen des Schiffchens, zwischen denen die gesuchte liegt, und wendet auf die jeweiligen Skalenablesungen ein einfaches Interpolationsverfahren an. Man kann dabei den Magnetstab in die natürliche oder in die umgekehrte Lage bringen. Im ersten Fall ergibt sich für $B$ (die Stellung der Alhidade des Schiffchens) der Wert $=Q$, im zweiten $=Q\pm180^0$. Die umgekehrte Lage bietet den Vorteil größerer Genauigkeit, weil einer kleinen Änderung von $B$ eine große Änderung der Skalenablesung entspricht. Die natürliche Lage ist dagegen etwas bequemer, da man das Schiffchen leichter so einstellen kann, dass die Grenzen der Skala eingehalten werden. Um mühsames Herumprobieren zu vermeiden, beginnt man daher am besten mit der natürlichen Lage und betrachtet das Ergebnis nur als Vorbereitung, damit man bei den Versuchen in umgekehrter Lage auf zwei nahe beieinanderliegende Teilstriche einstellen kann.

Das für $Q$ gefundene Ergebnis bezieht sich auf die Lage des magnetischen Meridians während oder zwischen den beiden letzten Versuchen. Ohne einen Hilfsapparat für Vergleichsbeobachtungen kann man keine zuverlässigere Bestimmung als diese schwankende Angabe erwarten. Steht aber ein Hilfsapparat zur Verfügung, so liefern gleichzeitige Stellungsbeobachtungen an ihm die Werte von $q$, die den beiden Beobachtungen entsprechen. Wendet man das obige Interpolationsverfahren auf die beiden Werte von $B-q$ an, so erhält man den Wert von $Q^0$ oder $Q^0\pm180^0$.

Schließlich beobachtet man die Schwingungsdauer sowohl in der natürlichen als auch in der umgekehrten Lage. Dazu stellt man die Alhidade des Schiffchens so genau wie möglich auf den Wert von $Q$ (beziehungsweise von $Q+180^0$) ein, der gerade zu Beginn der Schwingungsbeobachtungen gilt. Die Schwingungsdauer in der natürlichen Lage sei $t'$, in der umgekehrten $t''$. Kann man gleichzeitig Schwingungen am Hilfsmagnetometer beobachten, so erhält man daraus die entsprechenden Werte von $r$, die mit $r',r''$ bezeichnet seien. Will man auch die Veränderlichkeit von $S$ berücksichtigen, so seien $s',s''$ die entsprechenden Werte von $s$. Die kleinen Änderungen der Lage des magnetischen Meridians während der Schwingungsbeobachtungen
\SourceLeaf{421}werden in der Regel keinen merklichen Einfluss auf die Ergebnisse haben. Die beiden Gleichungen am Schluss von Artikel 2 lauten demnach
\[
\begin{aligned}
t't'&=\frac{s'S^0}{1+r'R^0}\\
t''t''&=\frac{s''S^0}{1-r''R^0}
\end{aligned}
\]
woraus durch Elimination folgt
\[
\begin{aligned}
R^0&=\frac{s't''t''-s''t't'}{r''s't''t''+r's''t't'}\\
S^0&=\frac{(r'+r'')t't't''t''}{r''s't''t''+r's''t't'}
\end{aligned}
\]

Nach der Bemerkung in Artikel 5 kann man ohne Weiteres $S$ als konstant betrachten oder $s'=s''=1$ setzen. Damit gehen die Formeln über in
\[
\begin{aligned}
R^0&=\frac{t''t''-t't'}{r''t''t''+r't't'}\\
S&=\frac{(r'+r'')t't't''t''}{r''t''t''+r't't'}
\end{aligned}
\]
Kann man aber keine Vergleichsbeobachtungen an einem Hilfsapparat heranziehen, so bleibt nur, unmittelbar
\[
\begin{aligned}
R&=\frac{t''t''-t't'}{t''t''+t't'}\\
S&=\frac{2t't't''t''}{t''t''+t't'}
\end{aligned}
\]
zu setzen. Der so gefundene Wert von $R$ ist offenbar nur eine Art Mittel zwischen den für die beiden Schwingungssätze geltenden Werten. $S$ bleibt mit einer kleinen Ungenauigkeit behaftet, die von der Ungleichheit dieser Sätze abhängt.

\Article{7}
Die allgemeinere Lösung unserer Aufgabe stützen wir auf gleichzeitige Beobachtungen der Schwingungsdauer und der Gleichgewichtsstellung des Bifilarmagnetometers bei zwei beliebigen verschiedenen Stellungen des Schiffchens. Wir bezeichnen die bestimmten Werte der Größen $A,B,p,Q,R,S,t$

für den ersten Beobachtungssatz mit $A',B',p',Q^0+q',r'R^0,s'S^0,t'$,

für den zweiten mit $A'',B'',p'',Q^0+q'',r''R^0,s''S^0,t''$.

\SourceLeaf{422}Setzt man diese Werte in die beiden Gleichungen (1) und (2) aus Artikel 2 ein, so ergeben sich vier Gleichungen, aus denen sich die unbekannten Größen $P,Q^0,R^0,S^0$ durch Elimination ableiten ließen. Durch Rechnen mit imaginären Größen gelangen wir viel leichter zum selben Ziel. Nach Formel (4) aus Artikel 2 gehen wir von den beiden Gleichungen aus
\[
\begin{aligned}
1&=\frac{s'S^0}{t't'}e^{i(P+p'-A')}-r'R^0e^{i(Q^0+q'-B')}\\
1&=\frac{s''S^0}{t''t''}e^{i(P+p''-A'')}-r''R^0e^{i(Q^0+q''-B'')}
\end{aligned}
\]
die, wenn wir zur Abkürzung
\[
\begin{aligned}
\frac{s'}{t't'}e^{i(p'-A')}&=a'\\
\frac{s''}{t''t''}e^{i(p''-A'')}&=a''\\
r'e^{i(q'-B')}&=b'\\
r''e^{i(q''-B'')}&=b''\\
S^0e^{iP}&=x\\
R^0e^{iQ^0}&=y
\end{aligned}
\]
setzen, übergehen in
\[
\begin{aligned}
1&=a'x-b'y\\
1&=a''x-b''y
\end{aligned}
\]
woraus man
\[
\begin{aligned}
x&=\frac{b''-b'}{a'b''-a''b'}=\frac{\dfrac{b''}{b'}-1}{\dfrac{b''}{b'}-\dfrac{a''}{a'}}.\frac{1}{a'}\\[1ex]
y&=\frac{a''-a'}{a'b''-a''b'}=\frac{\dfrac{a''}{a'}-1}{\dfrac{b''}{b'}-\dfrac{a''}{a'}}.\frac{1}{b'}
\end{aligned}
\]
erhält. Daraus ergeben sich die folgenden ausführlichen Rechenvorschriften. Man setze (I)
\SourceLeaf{423}
\[
\begin{aligned}
\frac{t't'}{s'}.\frac{s''}{t''t''}.\cos(A'-A''-p'+p'')&=\mathfrak{A}\\
\frac{t't'}{s'}.\frac{s''}{t''t''}.\sin(A'-A''-p'+p'')&=\mathfrak{A}_1\\
\frac{r''}{r'}.\cos(B'-B''-q'+q'')&=\mathfrak{B}\\
\frac{r''}{r'}.\sin(B'-B''-q'+q'')&=\mathfrak{B}_1
\end{aligned}
\]
sodass
\[
\begin{aligned}
\frac{a''}{a'}&=\mathfrak{A}+i\mathfrak{A}_1\\
\frac{b''}{b'}&=\mathfrak{B}+i\mathfrak{B}_1
\end{aligned}
\]
gilt. Ferner bestimme man die sechs Größen $u,U,v,V,w,W$ aus den Gleichungen (II)
\[
\begin{aligned}
\mathfrak{A}-1&=u\cos U\\
\mathfrak{A}_1&=u\sin U\\
\mathfrak{B}-1&=v\cos V\\
\mathfrak{B}_1&=v\sin V\\
\mathfrak{B}-\mathfrak{A}&=w\cos W\\
\mathfrak{B}_1-\mathfrak{A}_1&=w\sin W
\end{aligned}
\]
und zwar so, dass $u,v,w$ positiv sind. Dann gilt
\[
\begin{aligned}
\frac{b''}{b'}-1&=ve^{iV}\\
\frac{a''}{a'}-1&=ue^{iU}\\
\frac{b''}{b'}-\frac{a''}{a'}&=we^{iW}
\end{aligned}
\]
und folglich
\[
\begin{aligned}
x&=\frac{t't'}{s'}.\frac{v}{w}e^{i(V-W+A'-p')}\\
y&=\frac{u}{r'w}e^{i(U-W+B'-q')}
\end{aligned}
\]
woraus man leicht schließt, dass (III)
\SourceLeaf{424}
\[
\begin{aligned}
P&=V-W+A'-p'\\
Q^0&=U-W+B'-q'\\
R^0&=\frac{u}{r'w}\\
S^0&=\frac{t't'}{s'}.\frac{v}{w}
\end{aligned}
\]
Die vierzehn Formeln I, II, III enthalten die vollständige Lösung unserer Aufgabe in der einfachsten möglichen Form.

Noch ist zu bemerken, dass für $r'=r''$ (sei es, dass die Vergleichsbeobachtungen diese Gleichheit ergeben, oder dass wir mangels solcher Beobachtungen die Veränderlichkeit von $R$ während der beiden Beobachtungssätze nicht berücksichtigen können)
\[
\begin{aligned}
V&=\tfrac12(B'-B''-q'+q'')\pm90^0\\
v&=\pm2\sin\tfrac12(B'-B''-q'+q'')
\end{aligned}
\]
gilt. Dabei gelten die oberen oder unteren Zeichen, je nachdem, ob
\[
\sin\tfrac12(B'-B''-q'+q'')
\]
positiv oder negativ ist.

\Article{8}
Zur Erläuterung dieser Vorschriften fügen wir die vollständige Berechnung eines Beispiels bei. Die Rechnung ist mit siebenstelligen Logarithmen durchgeführt und damit viel genauer, als es für die praktische Anwendung nötig ist, bei der fünfstellige Logarithmen immer ausreichen.

Am 24. März 1841 wurde die Schwingungsdauer des Bifilarmagnetometers aus Beobachtungen bestimmt, die $1^{St}21'$ umfassten. Nach der gebotenen Reduktion auf unendlich kleine Schwingungen, die sich von selbst versteht, ergab sie sich zu $=28''89071=t'$. Die Stellung der Spiegelalhidade war $154^020'30''=A'$, die der Alhidade des Schiffchens $=27^040'25''=B'$. Im Mittel mehrerer gleichmäßig über diesen Zeitraum verteilter Bestimmungen betrug die Stellung $994.33$ Skalenteile. Da der Lotfaden der Skalenzahl $1000$ entspricht und ein Skalenteil $21''5835$ beträgt, ist $p'=-2'2''38$. Aus genau gleichzeitigen Beobachtungen ergab sich die Schwingungsdauer des Unifilarmagnetometers im magnetischen
\SourceLeaf{425}Observatorium zu $=20''72725$ und die mittlere Stellung zu $=881.80$ Skalenteilen. Als Normalstellung wurde der Mittelwert der täglichen Aufzeichnungen im Februar, $888.40$, gewählt (dem übrigens die absolute Deklination $18^011'54''$ entspricht). Da ein Skalenteil am Unifilarmagnetometer $21''3489$ beträgt, ergibt sich daraus $q'=-2'20''90$.

Der mittlere Thermometerstand (aus Aufzeichnungen unmittelbar vor Beginn und gleich nach dem Ende der Beobachtungen) betrug im Kasten des Bifilarmagnetometers $+6^096$, an der oberen Befestigung der Aufhängungsdrähte $+7^06$ und im Kasten des Unifilarmagnetometers $+7^045$, alles nach Réaumur.

Auf dieselbe Weise ergab sich für einen zweiten Beobachtungssatz am folgenden Tag
\[
\begin{aligned}
t''&=108''17\\
A''&=151^027'30''\\
B''&=185\;59\;35\\
p''&=-24'33''07\\
q''&=+2\;42.04
\end{aligned}
\]
die Schwingungsdauer des Unifilarmagnetometers zu $=20''73117$ und die Thermometerstände in derselben Reihenfolge wie oben zu $+6^036$, $+7^00$, $+7^01$.

Zur Berechnung des Temperatureinflusses setze ich
\[
\alpha=0.000016,\quad\beta=0.000024,\quad\gamma=\gamma'=0.000765
\]
den letzten Wert nach \textsc{Hansteen}, da eigene entscheidende Bestimmungen bisher noch fehlen. Daraus folgt nach den Formeln (1) und (3) von Artikel 5, wenn wir $20''72=\theta^0$ als Normalschwingungsdauer des Unifilarmagnetometers wählen,
\[
\begin{aligned}
\log r'&=-0.0001376\\
\log s'&=-0.0000043\\
\log r''&=-0.0002155\\
\log s''&=-0.0000044
\end{aligned}
\]
Nach den abgekürzten Formeln (2) und (4) derselben Stelle könnte man
\[
\begin{aligned}
\log s&=-\alpha k(c''-c)\\
\log r&=2\log\frac{\theta^0}{\theta}+(\gamma+2\alpha)k(c'-c)-\alpha k(c''-c)
\end{aligned}
\]
\SourceLeaf{426}setzen, wenn $k$ den Modul der Briggsschen Logarithmen bezeichnet, also mit den obigen Werten von $\alpha,\beta,\gamma$
\[
\begin{aligned}
\log s&=0.00000695(c''-c)\\
\log r&=2\log\frac{\theta^0}{\theta}+0.0003461(c'-c)-0.00000695(c''-c)
\end{aligned}
\]
woraus für unsere Beobachtungen folgt
\[
\begin{aligned}
\log r'&=-0.0001386\\
\log s'&=-0.0000044\\
\log r''&=-0.0002169\\
\log s''&=-0.0000044
\end{aligned}
\]
also kaum merklich von den obigen Werten verschieden.

Nach diesen Vorbereitungen sind die Hauptschritte der Rechnung selbst die folgenden
\[
\begin{aligned}
A'-A''-p'+p''&=+2^030'29''31\\
B'-B''-q'+q''&=-158\;14\;7.06\\
\log\frac{t't's''}{s't''t''}&=8.8524525\\
\log\frac{r''}{r'}&=9.9999221
\end{aligned}
\]
Hieraus nach I
\[
\begin{aligned}
\log\mathfrak{A}&=8.8520363\\
\log\mathfrak{A}_1&=7.4935432\\
\log\mathfrak{B}&=9.9678043\,\mathrm{n}\\
\log\mathfrak{B}_1&=9.5690569\,\mathrm{n}
\end{aligned}
\]
woraus ferner folgt
\[
\begin{aligned}
\log(\mathfrak{A}-1)&=9.9679562\,\mathrm{n}\\
\log(\mathfrak{B}-1)&=0.2852304\,\mathrm{n}\\
\log(\mathfrak{B}-\mathfrak{A})&=9.9998589\,\mathrm{n}\\
\log(\mathfrak{B}_1-\mathfrak{A}_1)&=9.5726915\,\mathrm{n}
\end{aligned}
\]
Danach ergeben die Formeln II
\SourceLeaf{427}
\[
\begin{aligned}
U&=179^048'28''15\\
V&=190\;52\;52.91\\
W&=200\;30\;14.79\\
\log u&=9.9679586\\
\log v&=0.2931101\\
\log w&=0.0282829
\end{aligned}
\]
und schließlich die Formeln III
\[
\begin{aligned}
P^0&=144^045'10''50\\
Q^0&=7\;0\;59.26\\
\log R^0&=9.9398133\\
\log S^0&=3.1863476
\end{aligned}
\]

\Article{9}
Die Aufgabe lässt sich noch weiter verallgemeinern, indem man \textit{vier verschiedene} Beobachtungssätze zugrunde legt, zwei für die Stellung und zwei für die Schwingungsdauer. Dabei lässt man die Voraussetzung fallen, dass die einen und die anderen jeweils denselben Werten von $B$ entsprechen. Das hat den Vorteil, dass man die Stellung des Bifilarmagnetometers nach dem in den \textit{Resultaten} für 1836. II. [siehe unten] beschriebenen Verfahren bei beinahe völlig beruhigtem Magnetstab beobachten kann. Dieser Vorteil verliert aber seinen Wert, da man für alle vier Sätze am Hilfsmagnetometer Schwingungsdauer und Stellung zugleich beobachten muss, also letztere doch aus Elongationen bestimmen muss. Diese Methode verursacht demnach offenbar doppelt so viel Arbeit wie die in Artikel 7, die außerdem den Vorteil einer sehr einfachen Rechnung hat. Die unmittelbare Bestimmung der Größen aus vier getrennten Beobachtungssätzen ist dagegen weit umständlicher. Ihre Entwicklung, die mathematisch nicht uninteressant ist, heben wir daher lieber für eine andere Stelle auf.

\Article{10}
Es ist noch zu beachten, dass, wenn man bei der Bestimmung der Stellung aus beobachteten Elongationen das in den [S. unten] \textit{Resultaten} 1836 II
\SourceLeaf{428}angegebene Verfahren ohne Weiteres anwendet, die unterschiedlichen Winkelwerte der Skalenteile einen Fehler verursachen. Er ist umso größer, je weiter die Stellung von der Mitte der Skala entfernt ist. Verlangt man also ganz genaue Ergebnisse, muss man das Verfahren auf die nach einer exakten Formel in Bogenteile umgerechneten Abstände der Elongationen von der Mitte der Skala anwenden, statt unmittelbar auf die bei den Elongationen abgelesenen Skalenteile. Liegt die Stellung nahe bei der Mitte, so ist dieser Fehler allerdings unerheblich. Man wird daher die Stellung des Spiegels oder den Wert von $A$ immer so wählen, dass die Stellung wenig von der Mitte abweicht, oder dass $p$ klein wird. Bei der ersten Bestimmung der Größen ist das freilich nur durch einen vorläufigen Versuch zu erreichen (ähnlich wie in Artikel 6). Kennt man aber die Größen $P,Q,R$ bereits annähernd, so wird man hierzu lieber eine Rechnung verwenden, der man am bequemsten die folgende, aus Artikel 2, Formel (1) oder (4) leicht abzuleitende Form gibt. Man bestimme einen Winkel $\psi$ durch die Formel
\[
\operatorname{tang}\psi=\frac{1-R}{1+R}.\operatorname{tang}\tfrac12(Q-B)=\operatorname{tang}\tfrac12\varphi^2.\operatorname{tang}\tfrac12(Q-B)
\]
und zwar so, dass $\psi$ im selben Quadranten gewählt wird wie $\tfrac12(Q-B)$, und setze dann
\[
A=\psi+P-\tfrac12(Q-B)
\]

\Article{11}
Nun bleibt noch zu zeigen, wie die Beobachtungen mit dem Bifilarmagnetometer in der Transversalstellung mit den Änderungen der Größen zusammenhängen.

Wir setzen voraus, dass sich die nach der Vorschrift in Artikel 3 bestimmte Transversalstellung auf die Normalwerte der Größen bezieht. Das Schiffchen ist also so eingestellt, dass im Ruhezustand die magnetische Achse des Magnetstabs mit dem magnetischen Normalmeridian einen rechten Winkel bildet, wenn sich die magnetische Richtungskraft zur statischen wie $R^0$ zu $1$ verhält. Der Spiegel ist dagegen so eingestellt, dass in jener Stellung das Bild des vom Lotfaden bezeichneten Skalenpunkts auf dem Fadenkreuz des Beobachtungsfernrohrs erscheint. Bezeichnen wir also die Schwingungsdauer unter diesen Umständen mit $t^0$, so gilt
\SourceLeaf{429}
\[
\begin{aligned}
\sin\varphi&=R^0\\
A&=P\mp\varphi\\
B&=Q^0\mp(90^0+\varphi)\\
\frac{S^0}{t^0t^0}&=\cos\varphi
\end{aligned}
\]
wobei sich die doppelten Vorzeichen auf die westliche beziehungsweise östliche Stellung des Nordpols des Magnetstabs beziehen. Behalten wir nun die Zeichen
\[
p,\quad Q=Q^0+q,\quad R=rR^0,\quad S=sS^0,\quad t
\]
in ihrer bisherigen allgemeinen Bedeutung bei, so ergeben die Formeln (1) und (6) des 2. Artikels
\[
\begin{aligned}
\sin(\varphi\pm p)&=rR^0\cos(p-q)\\
\frac{sS^0}{tt}\sin(\varphi\pm p)&=rR^0\cos(\varphi\pm q)
\end{aligned}
\]
oder
\[
\begin{aligned}
r&=\frac{\sin(\varphi\pm p)}{\sin\varphi\cos(p-q)}\quad .\ .\ .\ .\ .\ .\ .\ .\ .\quad(1)\\
\frac{st^0t^0}{tt}&=\frac{\cos(\varphi\pm q)}{\cos\varphi\cos(p-q)}\quad .\ .\ .\ .\ .\ .\ .\ .\ .\quad(2)
\end{aligned}
\]

\Article{12}
Die wichtigste Anwendung des Bifilarmagnetometers ist die Bestimmung der Änderungen der horizontalen Intensität. Mit diesen hängen die Änderungen von $R$ durch die oben (Artikel 5) angegebene Formel
\[
r=\frac{f^0g^0hMT}{fgh^0M^0T^0}
\]
zusammen. Dabei ist zu beachten, dass $T^0$ die anfangs gewählte Normalintensität bezeichnet und $M^0$ das auf die Normaltemperatur reduzierte magnetische Moment des Magnetstabs entsprechend seinem magnetischen Zustand zur Zeit der Bestimmung der Konstanten. Bezeichnen wir das ebenfalls auf die Normaltemperatur reduzierte magnetische Moment zu einer unbestimmten Zeit mit $\mathfrak{M}$ und setzen
\[
\frac{\mathfrak{M}}{M^0}=\lambda,\qquad\frac{T^0}{\lambda}=\mathfrak{T}
\]
so ist $\lambda$ unter den im 4. Artikel besprochenen Bedingungen ein Koeffizient, der für eine mäßig lange
\SourceLeaf{430}Zeit, etwa einen Tag, als konstant angesehen werden kann. Wie dieser zusammen mit $\mathfrak{M}$ nach und nach sehr langsam abnimmt, so nimmt $\mathfrak{T}$ allmählich zu und bezeichnet stets diejenige horizontale Intensität, bei der das Verhältnis der magnetischen zur statischen Richtungskraft bei Normaltemperatur gleich dem Verhältnis $R^0:1$ ist. Die obige Gleichung nimmt nun die Form
\[
T=\frac{\mathfrak{M}fgh^0}{Mf^0g^0h}.r\mathfrak{T}
\]
an. Der erste Faktor $\frac{\mathfrak{M}fgh^0}{Mf^0g^0h}$ hängt allein von der Temperatur ab und lässt sich (wenn wir die Bezeichnungen des 5. Artikels beibehalten) durch
\[
1+(\gamma+2\beta-\alpha)c+(\theta-\tfrac12\alpha)(c''-c)
\]
ausdrücken. Dagegen lässt sich $r$ nach Formel (1) des vorigen Artikels für jeden Augenblick bestimmen, indem man gleichzeitige Standbeobachtungen mit dem Bifilarmagnetometer in der Transversalstellung für $p$ und mit dem Unifilarmagnetometer für $q$ miteinander verbindet. Auf diese Weise lassen sich also die Änderungen der Intensität in den kürzesten Zeitabständen mit größter Genauigkeit verfolgen, solange es nur darauf ankommt, die veränderten Intensitäten über einen mäßig langen Zeitraum \textit{untereinander} zu vergleichen, etwa während eines vierundzwanzigstündigen Beobachtungstermins oder während der Zeit, die für eine absolute Intensitätsbestimmung mit dem Unifilarmagnetometer erforderlich ist. Verwendet man bei einer solchen absoluten Intensitätsbestimmung die gleichzeitigen Beobachtungen mit dem Bifilarmagnetometer, um die einzelnen Operationen auf dieselbe Normalintensität zu reduzieren (vgl. \textit{Intensitas vis magneticae}, Artikel 10 und 22), so erhält man zugleich den für diese Zeit gültigen Wert von $\mathfrak{T}$ in absolutem Maß. Dies ist auch an sich vorteilhafter als der dort empfohlene Gebrauch eines zweiten Unifilarmagnetometers. Wiederholt man solche absoluten Bestimmungen von Zeit zu Zeit, so bleibt man über etwaige allmähliche Änderungen von $\mathfrak{T}$ laufend unterrichtet. Man kann sie für die Zwischenzeit durch Interpolation in geometrischer Folge ohne merklichen Fehler ansetzen und so sämtliche Änderungen der Intensität in all ihren Schwankungen in absolutem Maß angeben. Im Übrigen versteht es sich von selbst, dass man, wenn nach einem längeren Zeitraum infolge der Säkularänderungen der magnetischen Deklination und der horizontalen In%
\SourceLeaf{431}tensität oder einer beträchtlichen Schwächung des Stabmagnetismus $p$ und $q$ nicht mehr innerhalb mäßiger Grenzen bleiben (Fälle großer außerordentlicher Anomalien braucht man dabei nicht mitzuzählen), die Stellung des Schiffchens und des Spiegels und, wenn es ratsam erscheint, auch den Abstand der Aufhängungsdrähte zweckmäßig ändert und so eine neue Beobachtungsreihe mit geänderten Größen beginnt.

\Article{13}
Solange $p$ und $q$ klein sind, kann man für alle Zwecke, die keine höchste Genauigkeit erfordern, anstelle der exakten Formel (1) eine verkürzte verwenden. In ihr fällt $q$ ganz weg, sodass man gleichzeitige Beobachtungen mit dem Unifilarmagnetometer überhaupt nicht braucht. Dies gilt insbesondere für die gewöhnlichen Terminsbeobachtungen. Anstelle jener Formel kann man nämlich setzen
\[
r=1\pm\operatorname{cotang}\varphi.\operatorname{tang}p\quad\text{oder auch}\quad r=1\pm\tfrac12\operatorname{cotang}\varphi.\operatorname{tang}2p
\]
Wenn nun $n$ den Unterschied zwischen dem abgelesenen Skalenteil und demjenigen bezeichnet, auf den sich der Lotfaden bezieht, und $d$ die horizontale Entfernung der Spiegelmitte von diesem letzteren Punkt, gemessen in Skalenteilen, dann gilt
\[
\operatorname{tang}2p=\frac{n}{d}
\]
und die Formel geht über in
\[
r=1\pm\frac{n}{2\operatorname{tang}\varphi.d}
\]
Zugleich erhält man mit hinreichender Genauigkeit
\[
T=\mathfrak{T}(1\pm\frac{n}{2\operatorname{tang}\varphi.d}+(\gamma+2\beta-\alpha)c)
\]
wenn man das geringfügige Glied $(\beta-\tfrac12\alpha)(c''-c)$ weglässt. Beim hiesigen Apparat ist $d=4778{,}3$ Millimeter, und aus den Ergebnissen der in Artikel 8 als Beispiel ausgeführten Rechnung ergibt sich $\varphi=60^031'37''9$, also $2d\operatorname{tang}\varphi=16910$. Mit den dort verwendeten Werten von $\alpha,\beta,\gamma$ erhält man also
\[
T=(1+\frac{n+13{,}65c}{16910})\mathfrak{T}
\]
\SourceLeaf{432}wenn sich der Nordpol des Magnetstabs auf der Westseite befindet, und
\[
T=(1-\frac{n-13{,}65c}{16910})\mathfrak{T}
\]
wenn er sich auf der Ostseite befindet.

Im Übrigen versteht es sich, dass man die Temperatur bei den Terminsbeobachtungen ohne Weiteres ganz außer Acht lassen kann, solange es nur darum geht, die Gestalt der einzelnen, in kurzen Zeitabständen wechselnden Anomalien zu erkennen.

\Article{14}
In der Transversalstellung des Bifilarmagnetometers wirken sich die Änderungen der Intensität in ihrer ganzen Stärke, die der Deklination dagegen kaum merklich auf die Stellung aus. Gerade umgekehrt haben die letzteren Änderungen den größten, die ersteren dagegen nur einen äußerst geringen Einfluss auf die Schwingungsdauer. Sofern $p,q$ und die Abweichung der Größe $S$ vom Normalwert klein sind, kann man ohne erheblichen Fehler anstelle von Formel (2) in Artikel 11
\[
\frac{t^0t^0}{tt}=1\mp q\operatorname{tang}\varphi\quad\text{oder auch}\quad t=t^0(1\pm\tfrac12q\operatorname{tang}\varphi)
\]
setzen, wenn $q$ in Teilen des Halbmessers ausgedrückt ist, und folglich
\[
t=t^0\pm\frac{t^0\operatorname{tang}\varphi}{412530}.q
\]
wenn es in Bogensekunden ausgedrückt ist. Aus den Ergebnissen des oben berechneten Beispiels folgt $t^0=55{,}871$ Zeitsekunden, sodass man in Bogensekunden
\[
q=\pm(t-55{,}871)4172''8
\]
erhält. Die exakte Umformung der Formel (2) zur Berechnung von $q$ lautet
\[
\pm\operatorname{tang}q=\frac{tt-st^0t^0\cos p}{tt\operatorname{tang}\varphi\pm st^0t^0\sin p}
\]
Im Übrigen versteht es sich, dass sich auf diese Weise durch Schwin\-gungsbeobachtungen kein für einen bestimmten Augenblick gültiger Wert von $q$ bestimmen lässt, sondern nur der Mittelwert über die Dauer dieser Beobachtungen.

